\documentclass[a4paper]{article}

\usepackage[utf8]{inputenc}
\usepackage[spanish]{babel}
\usepackage{amsmath,amssymb,amsfonts}
\usepackage{graphicx}
\usepackage{enumitem}
\usepackage{verbatim}
\usepackage{color}

% Definir los entornos necesarios para exams
\newenvironment{question}{\item}{}
\newenvironment{solution}{\comment}{\endcomment}
\newenvironment{answerlist}{\renewcommand{\labelenumii}{(\alph{enumii})}\begin{enumerate}}{\end{enumerate}}

%% comandos para exams
\newcommand{\exname}[1]{\def\@exname{#1}}
\newcommand{\extype}[1]{\def\@extype{#1}}
\newcommand{\exsolution}[1]{\def\@exsolution{#1}}
\newcommand{\exshuffle}[1]{\def\@exshuffle{#1}}
\newcommand{\exsection}[1]{\def\@exsection{#1}}

\usepackage{Sweave}
\begin{document}
\input{costo_boleta_interpretacion_representacion_n2_v1-concordance}

\begin{enumerate}




\begin{question}

En un circo, el presupuesto total disponible para entretenimiento es de $80000$. Este presupuesto se distribuye entre el costo de boleta y el costo de dulces, de tal manera que cuando aumenta el gasto en uno, disminuye proporcionalmente el gasto en el otro.

La siguiente tabla muestra la relacion entre estas dos variables:

% Crear tabla de datos simple
\begin{center}
\begin{tabular}{|c|c|}
\hline
\textbf{Costo de boleta (x)} & \textbf{Costo de dulces (y)} \\
\hline
0 & 80000 \\
\hline
8000 & 72000 \\
\hline
16000 & 64000 \\
\hline
24000 & 56000 \\
\hline
32000 & 48000 \\
\hline
40000 & 40000 \\
\hline
48000 & 32000 \\
\hline
56000 & 24000 \\
\hline
\end{tabular}
\end{center}

La grafica que representa esta funcion es:

\begin{answerlist}
\item
\includegraphics{costo_boleta_interpretacion_representacion_n2_v1-002}

\item
\includegraphics{costo_boleta_interpretacion_representacion_n2_v1-003}

\item
\includegraphics{costo_boleta_interpretacion_representacion_n2_v1-004}

\item
\includegraphics{costo_boleta_interpretacion_representacion_n2_v1-005}
\end{answerlist}

\end{question}

\begin{solution}

Para resolver este problema, debemos analizar la relacion entre las dos variables presentadas en la tabla.

Observando los datos:
- Cuando el boleta cuesta $0$, el dulces cuesta $80000$
- Cuando el boleta cuesta $56000$, el dulces cuesta $24000$

Esto muestra una relacion lineal decreciente: a medida que aumenta el costo de boleta, disminuye el costo de dulces, manteniendo un presupuesto total constante de $80000$.

\begin{answerlist}
  \item Falso. La grafica A muestra una relacion constante, no la relacion inversa esperada entre los costos.
  \item Verdadero. La grafica B representa correctamente la relacion lineal decreciente: a mayor costo de uno, menor costo del otro.
  \item Falso. La grafica C muestra una funcion escalonada, no una relacion continua como la del presupuesto fijo.
  \item Falso. La grafica D muestra puntos dispersos sin conexion, no una funcion matematica definida.
\end{answerlist}
\end{solution}

%% META-INFORMATION
\exname{Costo Boleta Interpretacion Representacion}
\extype{schoice}
\exsolution{0100}
\exshuffle{TRUE}
\exsection{Interpretacion y Representacion}

\end{enumerate}
\end{document}
